% Ensayo de Ledesma

En su artículo "The value of our personal data in the Big Data and the Internet of all Things Era", la autora Anahiby Becerril aborda la transformación socioeconómica provocada por la cuarta revolución industrial y una economía digital basada en la información. La relevancia del texto radica en exponer cómo nuestras interacciones con servicios digitales ''gratuitos'' nos hace ceder de forma casi involuntaria nuestra privacidad, muchas veces sin ser conscientes de valor que estamos generando. Esta dinámica resulta en una recolección masiva de datos sin un consentimiento informado o comprendido de las partes involucradas, lo que resulta en una pérdida de control sobre la información. Esta problemática planteada resulta muy relevante en el panorama actual y la postura de la autora me parece muy acertada, los usuarios deberían tener poder sobre el uso de sus datos.

La autora presenta la idea de que la economía digital depende fundamentalmente de los datos recolectados, pues empresas como Google, Facebook o Amazon obtienen grandes beneficios gracias a la recolección masiva de datos, convirtiendo a los usuarios en activos intangibles, que deben ceder, a veces de manera inconsciente, su información como moneda de cambio para acceder a los servicios ''gratuitos'' ofrecidos por estas empresas. Esto genera una asimetría de poder entre las empresas tecnológicas y los gobiernos que concentran dicha información, pues los individuos que generan los datos, carecen de control sobre ellos. 

Asimismo, se sostiene que el modelo económico actual perpertúa una explotación desproporcionada. Las empresas tecnológicas obtienen ganancias multimillonarias mediante el análisis de estos datos, mientras que los usuarios únicamente reciben el acceso a la plataforma, asumiendo todos los riesgos en materia de seguridad. Ejemplo de esto pueden ser las redes sociales, pues basan la totalidad de su valuación en la cantidad de usuarios activos y el volumen de datos que estos entregan a diario. Ejemplos como el caso Cambridge Analytica, donde millones de perfiles de Facebook fueron explotados para manipular campañas políticas, muestran los riesgos de no tener control.

Es por esto que el objetivo fundamental de la autora es informar al lector sobre la urgencia de reevaluar la importancia de los datos en la economía digital, tomando conciencia de las implicaciones de su recolección y promoviendo el empoderamiento mediante soluciones tecnológicas como los almacenes de datos personales (\textit{Personal Data Stores}) y normativas como la \textit{GDPR}. De esta forma, se promueve una economía digital más ética, justa y sostenible en la que la tecnología se desarrolla sin sacrificar los derechos de las personas.

Relacionando la problemática con la materia de Fundamentos de Bases de Datos, tenemos que las bases de datos no son simplemente estructuras de almacenamiento, sino que fungen como el núcleo donde se organizan, procesan e integran estos volúmenes masivos de información para ser convertidos en información útil. Es así que todo esto cobra un sentido ético y profesional, pues hay que tomar en cuenta que cada reigstro en una tabla corresponde a un atributo o comportamiento de la vida de una persona real.

Para concluir, reafirmo mi postura inicial, el control de los datos debe estar en manos de los usuarios, pues solo así se podrá equilibrar la relación entre individuos y empresas tecnológicas. El artículo muestra que los datos personales constituyen un activo de valor incalculable que no debe ser cedido sin consciencia, requiriendo transparencia y control normativo. La aportación más relevante del artículo me parece que es la necesidad de pasar de una postura pastiva como consumidores a un rol activo, reconociendo que la soberanía digital exige un conocimiento técnico y una regulación adecuada. Finalmente, este panorama plantea preguntas interesantes tales como ¿Cómo garantizar la privacidad de los datos ante el acelerado crecimiento del \textit{IoT} sin detener la innovación tecnológica? ¿De qué manera los profesionales del área de datos asumirán el compromiso ético de proteger la dignidad de las personas a través de sus datos?